\chapter{Outils de développement}

\section{Environnement intégré et langages}

\begin{itemize}[leftmargin=*]
  \item \textbf{Édition :} Visual Studio Code, Android Studio ou environnement équivalent pour Flutter.
  \item \textbf{Langages :} JavaScript (Node.js), Dart (Flutter), TypeScript (admin web), SQL (MySQL).
\end{itemize}

\section{Gestion de version et collaboration}

\begin{itemize}[leftmargin=*]
  \item \textbf{Git} pour le suivi des versions ;
  \item hébergement sur \textbf{GitHub} (branche principale, pull requests, workflow CI).
\end{itemize}

\section{Outils de build et d'exécution}

\begin{itemize}[leftmargin=*]
  \item \textbf{npm} pour le backend et l'admin ;
  \item \textbf{Flutter SDK} (canal stable) pour le mobile ;
  \item \textbf{Vite} pour le bundling de l'interface administrateur.
\end{itemize}

\section{Qualité et supervision}

\begin{itemize}[leftmargin=*]
  \item analyse statique Flutter (\texttt{flutter analyze}) ;
  \item tests backend (Mocha) et pipeline CI ;
  \item \textbf{Sentry} pour le suivi d'erreurs côté API et mobile (configuration par environnement).
\end{itemize}

\section{Rédaction du rapport}

\begin{itemize}[leftmargin=*]
  \item \textbf{\LaTeX} (moteur \texttt{pdflatex}) pour la composition du présent document ;
  \item \textbf{PlantUML} pour les diagrammes UML exportés en PNG dans \texttt{pfe/figures/}.
\end{itemize}
